% 第9章 从基态到激发
\chapter{从基态到激发}

前几章的重点主要在 Heisenberg 模型的基态。知道精确基态——或者退而求其次，一个好的变分态——使我们能够计算 $T = 0$ 的等时自旋关联。然而实验测量的是有限频率、有限温度下的动力学响应。由线性响应理论（见附录 B）显然可知，动力学的关联依赖于激发态与激发能。

本章将用基态关联构造某些近似低激发。这一途径称为单模近似（single mode approximation，SMA）\footnote{SMA 最初由 Bijl 与 Feynman 用于确定超流 ${}^{4}$He 的声子--旋子色散曲线。}。

Heisenberg 反铁磁体为展示具有高度关联基态的体系的 SMA 提供了机会。

自旋激发与关联的信息包含在动力学结构因子 (B.15) 中：

\begin{align*}
S ( \mathbf {q}, \omega ) & = \mathcal {N} ^ {- 1} \int_ {- \infty} ^ {\infty} d t \, e ^ {i \omega t} \sum_ {ij} e ^ {i \mathbf {q} ( \mathbf {X} _ {i} - \mathbf {X} _ {j} )} \langle S _ {i} ^ {z} ( t ) S _ {j} ^ {z} ( 0 ) \rangle\\
& = \frac {2 \pi}{\mathcal {N} Z} \sum_ {\alpha , \beta} e ^ {- E _ {\alpha} / T} \langle \alpha | S _ {\mathbf {q}} ^ {z} | \beta \rangle \langle \beta | S _ {- \mathbf {q}} ^ {z} | \alpha \rangle \, \delta ( \omega + E _ {\alpha} - E _ {\beta} ) .
\tag{9.1}
\end{align*}

等时关联函数为

\begin{align*}
S ( \mathbf {q} ) & = \int_ {- \infty} ^ {+ \infty} \frac {d \omega}{2 \pi} \, S ( \mathbf {q}, \omega )\\
& = \mathcal {N} ^ {- 1} \langle S _ {\mathbf {q}} ^ {z} S _ {- \mathbf {q}} ^ {z} \rangle ,
\tag{9.2}
\end{align*}

它只依赖基态波函数。另一个有用的函数是“双对易子”关联函数——先前在 Mermin--Wagner 定理的证明中遇到过（见 (6.22)）：

\begin{align*}
F ( \mathbf {q} ) & = \mathcal {N} ^ {- 1} \left\langle \left[ S _ {- \mathbf {q}} ^ {z}, [ \mathcal {H}, S _ {\mathbf {q}} ^ {z} ] \right] \right\rangle\\
& = \int_ {- \infty} ^ {+ \infty} \frac {d \omega}{2 \pi} \, \omega \, S ( \mathbf {q}, \omega ).
\tag{9.3}
\end{align*}

（这里与 (6.22) 中自旋方向的不同选择，对转动不变的哈密顿量无关紧要。）$F(\mathbf{q})$ 虽是静态基态期望值，却通过 $\mathcal{H}$ 包含动力学信息。

方程 (9.2) 与 (9.3) 表明，$S(\mathbf{q})$ 与 $F(\mathbf{q})$ 分别描述 $S(\mathbf{q},\omega)$ 的平均值与第一频率矩。于是，动量 $\mathbf{q}$ 处自旋激发的特征频率定义为

\begin{equation*}
\bar {\omega} _ {\mathbf {q}} \equiv \frac {F ( \mathbf {q} )}{S ( \mathbf {q} )}.
\tag{9.4}
\end{equation*}

\section{单模近似}

记基态为 $|0\rangle$。由平移不变性，所有激发都可用格子动量 $\mathbf{q}$ 标记。“单模”态构造为

\begin{equation*}
\left| \mathbf {q} \right\rangle = S _ {\mathbf {q}} ^ {z} \left| 0 \right\rangle .
\tag{9.5}
\end{equation*}

若 $T=0$ 的 $S(\mathbf{q},\omega)$ 在 $\omega=\bar{\omega}_{q}$ 附近有尖锐峰，即

\begin{equation*}
S ( \mathbf {q}, \omega ) \approx 2 \pi S ( \mathbf {q} ) \, \delta ( \omega - \bar {\omega} _ {\mathbf {q}} ) ,
\tag{9.6}
\end{equation*}

则单模态近似于体系的一个真实激发。由 (9.1) 第二行，若

\begin{equation*}
( \mathcal {H} - E _ {0} ) \left| \mathbf {q} \right\rangle = \bar {\omega} _ {\mathbf {q}} \left| \mathbf {q} \right\rangle ,
\tag{9.7}
\end{equation*}

则 (9.6) 是精确的。

一般而言，即使 (9.7) 严格成立，单模态也只是全部激发中通过 $S_{q}^{z}$ 与基态相连的一个小子集。单模态的数目随体系增大而增多。若它们的乘积也是近似激发，则可以把它们视为弱相互作用元激发。于是其能量近似地决定模型的低频响应与低温热力学。

6.3 节我们看到，基态与最低激发态之间在热力学极限下不消失的能隙 $\Delta$ 蕴含 $T = 0$ 无长程序。此时由 (9.1)，结构因子在能隙频率以下为零：

\begin{equation*}
S ( \mathbf {q}, \omega ) = 0 \, , \quad \omega \in ( 0, \Delta ) ,
\tag{9.8}
\end{equation*}

而由 (9.2)--(9.4)，单模能量 $\bar{\omega}_{\mathbf{q}}$ 是动量 $\mathbf{q}$ 处最低激发能的上界：

\begin{equation*}
\min_{\alpha} E _ {\alpha} ( \mathbf {q} ) \leq \bar {\omega} _ {\mathbf {q}}.
\tag{9.9}
\end{equation*}

遗憾的是，即使单模谱有能隙

\begin{equation*}
\Delta^ {\mathrm{SMA}} \leq \bar {\omega} _ {\mathbf {q}} \, , \quad \forall \mathbf {q} ,
\tag{9.10}
\end{equation*}

也不蕴含精确本征能量 $E_{\alpha}$ 有真正的能隙。

\section{Goldstone 模}

不等式 (9.9) 不能证明能隙的存在，却能证明无能隙。这里我们将用它证明 Heisenberg 模型 Goldstone 定理的格子版本。Goldstone 定理的实质是：对具有短程相互作用（且无规范场）的哈密顿量，自发破缺的连续对称蕴含低能激发的存在，即 Goldstone 模。若基态动量为 $\bar{q}$，则 Goldstone 模的能量在 $q \rightarrow \bar{q}$ 时消失。例如，铁磁体 $\bar{q} = 0$，立方格子上的 Néel 反铁磁体 $\bar{q} = \vec{\pi}$。在相对论场论中，Goldstone 模代表无质量粒子。在凝聚态物理中，Goldstone 模出现在许多体系：如破缺平移对称的固体中的声学声子；$O(n)$ Heisenberg 模型（$n > 1$，$n$ 为自旋维数）中的自旋波。

Goldstone 定理的证明由 Lange 详细给出（见文献注释）。这里给出一个较短的版本\footnote{作者感谢 Daniel Arovas 提供此证明。}，它利用 SMA 上界 (9.9)。

考虑满足 Mermin--Wagner 定理条件（见 6.2 节）的短程 Heisenberg 哈密顿量：

\begin{equation*}
\mathcal {H} = \frac {1}{2} \sum_ {ij} J _ {ij} \, \mathbf {S} _ {i} \cdot \mathbf {S} _ {j} ,
\tag{9.11}
\end{equation*}

其中

\begin{equation*}
\bar {J} = \frac {1}{2 \mathcal {N}} \sum_ {ij} \left| J _ {ij} \right| \left| \mathbf {x} _ {i} - \mathbf {x} _ {j} \right| ^ {2} < \infty .
\tag{9.12}
\end{equation*}

\medskip
\noindent\textbf{Goldstone 定理 9.1}\quad 若自旋关联在某波矢 $\bar{\mathbf{q}}$ 发散：

\begin{equation*}
\lim _ {\mathbf {q} \to \bar {\mathbf {q}}} S ( \bar {\mathbf {q}} ) \to \infty ,
\tag{9.13}
\end{equation*}

则存在以动量 $\mathbf{q}$ 标记的 Goldstone 模，其能量 $E(\mathbf{q})$ 在 $\bar{q}$ 处消失：

\begin{equation*}
\lim _ {\mathbf {q} \to \bar {\mathbf {q}}} E ( \mathbf {q} ) = 0.
\tag{9.14}
\end{equation*}
\medskip

用 $h = 0$ 时 (6.26) 给出的 $F(\mathbf{q})$ 的上界：

\begin{equation*}
F ( \mathbf {q} ) \leq 2 S ( S + 1 ) \bar {J} \, | \mathbf {q} | ^ {2} ,
\tag{9.15}
\end{equation*}

结合 (9.4) 与 (9.15) 得

\begin{equation*}
\bar {\omega} _ {\bar {\mathbf {q}}} \leq \frac {2 S ( S + 1 ) \bar {J} | \bar {\mathbf {q}} | ^ {2}}{S ( \mathbf {q} )}.
\tag{9.16}
\end{equation*}

于是由 (9.9) 与 (9.13)，方程 (9.14) 得证。证毕。

\medskip
\noindent\textbf{推论 9.2}\quad 若基态有破缺对称，则如定理 9.2，存在 Goldstone 模。
\medskip

这直接由 (6.6) 得出：破缺对称蕴含真长程序，且 $\mathcal{N} \to \infty$ 时 $S(\mathbf{q})$ 发散。证毕。

重要的是，Goldstone 定理的逆命题不成立，即无能隙激发的存在不蕴含真长程序。著名的反例是 Fermi 气体：它有无能隙的粒子--空穴激发，但没有长程序\footnote{Fermi 液体（若存在的话）同样无能隙而无破缺对称。}。另外，一维 $S = \frac{1}{2}$ Heisenberg 反铁磁体的激发在 $q = 0$、$\pi$ 处消失：

\begin{equation*}
\omega_ {q} \propto | \sin q | ,
\tag{9.17}
\end{equation*}

但其关联在长距离上按代数方式衰减。

\section{Haldane 能隙与 SMA}

8.3 节已证明 AKLT 模型的基态是价键固体。$S = 1$ 的一维模型为

\begin{equation*}
\mathcal {H} ^ {\mathrm{AKLT}} = K \sum_{\langle ij \rangle} \left[ \mathbf {S} _ {i} \cdot \mathbf {S} _ {j} + \frac {1}{3} \left( \mathbf {S} _ {i} \cdot \mathbf {S} _ {j} \right) ^ {2} + \frac {2}{3} \right].
\tag{9.18}
\end{equation*}

$S = 1$ 的 VBS 态的关联函数已由 (8.30) 给出。取其 Fourier 变换：

\begin{equation*}
S ( q ) = \frac {2 ( 1 - \cos q )}{( 5 + 3 \cos q )}.
\tag{9.19}
\end{equation*}

对价键波函数同样可以计算 $F(q)$，结果\footnote{见 Arovas、Auerbach 与 Haldane，第 8 章文献注释。}为

\begin{equation*}
F ( q ) = \frac {10 K}{27} ( 1 - \cos q ).
\tag{9.20}
\end{equation*}

由 (9.4)，单模能量为

\begin{equation*}
\bar {\omega} _ {\mathbf {q}} = \frac {5 K}{27} ( 5 + 3 \cos q ) \geq 0. 3 7 0 \, K.
\tag{9.21}
\end{equation*}

如前所述，虽然单模谱有大小 $0.370K$ 的能隙，这并不蕴含精确谱中确实存在能隙。然而对 $S = 1$ AKLT 模型的数值模拟支持存在大小 $\Delta = 0.350K$ 的能隙\footnote{这是 F. D. M. Haldane 未发表的结果。}。对标准 Heisenberg 模型的同一模拟发现能隙大小为 $\Delta = 0.325K$。由于 AKLT 模型的 SMA 谱有能隙，这暗示附加的双二次项 $\frac{K}{3}(\mathbf{S}_{i} \cdot \mathbf{S}_{j})^{2}$ 并不剧烈改变一维 $S = 1$ Heisenberg 模型的基态关联与低激发\footnote{双二次微扰对关联的影响见 Arovas 与 Girvin（第 8 章文献注释）。}。

这个通常称为 Haldane 能隙的能隙在热力学极限（$\mathcal{N} \to \infty$）下存留。它是低维整数自旋量子反铁磁体的一般特征，由 Haldane 的连续理论在第 12、15 章解释。

\begin{chapbib}
\item 超流 ${}^{4}$He 的单模近似描述见
  R. P. Feynman, \textit{Statistical Mechanics} (Benjamin, New York, 1972), 第十一章。
\item SMA 被用于导出分数量子霍尔体系的磁声子与磁旋子激发：
  S. M. Girvin, A. H. Macdonald, and P. M. Platzman, Phys. Rev. B \textbf{33}, 2481 (1986)。
\item 铁磁体的 Goldstone 定理证明于
  R. V. Lange, Phys. Rev. \textbf{146}, 301 (1966)。
\item 推荐关于自发对称破缺与 Goldstone 模的讲义：
  S. Coleman, \textit{Aspects of Symmetry} (Cambridge, 1985), 第五章。
\end{chapbib}
